{\footnotesize\setlength{\tabcolsep}{4pt}\renewcommand{\arraystretch}{1.0}
\begin{tabular}{@{}>{\raggedright\arraybackslash}p{0.24\linewidth}>{\raggedright\arraybackslash}p{0.72\linewidth}@{}}
\texttt{chunk\_id} & labour\_952\_s340 \\
\texttt{provision\_id} & labour\_952\_s340 \\
\texttt{act\_id} & labour\_952 \\
\texttt{act\_title\_bn} & \bng{বাংলাদেশ শ্রম আইন, ২০০৬} \\
\texttt{act\_year} & 2006 \\
\texttt{provision\_kind} & section \\
\texttt{provision\_no\_bn} & \bng{৩৪০} \\
\texttt{provision\_no\_ascii} & 340 \\
\texttt{provision\_title\_bn} & \bng{নিয়োগ সম্পর্ক অনুমান} \\
\texttt{domain} & labour \\
\texttt{risk\_tier} & medium \\
\texttt{n\_sub} & 1 \\
\texttt{source\_url} & http://bdlaws.minlaw.gov.bd/act-print-952.html \\
\texttt{source\_dataset} & blad \\
\texttt{crawl\_date} & 2025-07-19 \\
\texttt{n\_words} & 42 \\
\texttt{language} & bengali \\
\texttt{text\_raw} & \bng{কোন ব্যক্তিকে কোন কারখানায় উহার কাজ চলাকালে, আহার বা বিশ্রামের জন্য বিরতির সময় ব্যতীত অন্য কোন সময়ে অথবা উহার কোন যন্ত্রপাতি চালু থাকা কালে পাওয়া যায় তাহা হইলে, ভিন্নরূপ প্রমাণীত না হওয়া পর্যন্ত, তিনি উক্ত সময়ে কারখানায় নিযুক্ত ছিলেন বলিয়া গণ্য হইবেন৷} \\
\end{tabular}}
